\chapter*{Preface}
\phantomsection
\addcontentsline{toc}{chapter}{Preface}
\markboth{Preface}{Preface}

\section*{Systematic trading and investing}
\phantomsection
\addcontentsline{toc}{section}{Systematic trading and investing}

I am very bad at making financial decisions. Like most people I find it difficult to
manage my investments without becoming emotional and behaving irrationally. This is
deeply irritating as I consider myself to be very knowledgeable about finance. I've
voraciously read the academic literature, done my own detailed research, spent 20 years
investing my own money and nearly a decade managing funds for large institutions.

So in theory I know what I'm doing. In practice when faced with a decision to buy or sell
a stock things go wrong. Fear and greed wash through my mind, clouding my judgment.
Even if I've spent weeks researching a company it's still hard to click the trade button on
my broker's website. I have to stop myself buying or selling on a whim, based on nothing
more than random newspaper articles or an anonymous blogger's opinion. But then, like
you, I'm only human.

Fortunately there is a solution. The answer is to fully, or partly, systematise your
financial decision making. Creating a trading system removes the emotion and makes it
easier to commit to a consistent strategy.

I spent many years managing a large portfolio of trading strategies for a systematic hedge
fund. Unfortunately I didn't have the opportunity to develop and trade systems to look
after my personal portfolio. But after leaving the industry I've been able to make my own
trading process entirely systematic, resulting in significantly better performance.

There are many authors and websites offering trading systems. But many of these
``systems'' require subjective interpretation, so they are not actually systematic. Some are
even downright dangerous, trading too quickly and expensively, and in excessive size.
They present you with a single ``one size fits all'' system which won't suit everybody. I
will explain how to develop your own trading system, for your own needs, and which
should not be excessively dangerous or costly to operate.

I've found systematic investing to be more profitable, and to require less time and effort.
This book should help you to reap similar benefits.

\section*{Who should read this book}
\phantomsection
\addcontentsline{toc}{section}{Who should read this book}

This book is intended for everyone who wishes to systematise their financial decision
making, either completely or to some degree.

Most people would describe themselves as traders or investors, although there are no
consistently accepted definitions of either group. What I have to say is applicable to both
kinds of readers so I use the terms trading and investing interchangeably; the use of one
usually implies the other is included.

This book will be useful to amateurs --- individual investors trading their own money ---
and to market professionals who invest on behalf of others. The term ``amateur'' is not
intended to be patronising. It means only that you are not getting paid to manage money
and is no reflection on your level of skill.

This is not intended to be a parochial book solely for UK or US investors and I use
examples from a range of countries. Many books are written specifically for particular
asset markets. My aim is to provide a general framework that will suit traders of every
asset. I use specific examples from the equity, bond, foreign exchange and commodity
markets. These are traded with spread bets, exchange traded funds\footnote{All terms in
bold are defined in the glossary.} and futures. But I do not explain the mechanics of
trading in detail. If you are not familiar with a particular market you should consult other
books or websites before designing your trading system.

It might surprise you, but this book will also be useful for those who are sceptical of
computers entirely replacing human judgment. This is because there are several parts to a
complete trading system. Trading rules provide a prediction on whether something will go
up or down in price. These can be purely systematic, or based on human discretion. But
it is equally important to have a good framework of position and risk management. I
believe that a systematic framework should be used by all traders and investors for
position and risk management, even if the adoption of fully systematic rules is not
desirable.

If you can beat simple rules when it comes to predicting prices, I show you how to use
your opinions in a systematic framework to make the best use of your talents.
Alternatively you might feel it is unlikely that anyone, man or machine, can predict the
markets. In this case, the same framework can be used to construct the best portfolio
consistent with that pessimistic view.

\section*{Three examples}

Throughout this book I focus on three typical groups of systematic traders and investors.
Don't panic if you don't fit neatly into any of these categories. I've chosen them because
between them they illustrate the most important issues which face all potential systematic
traders and investors.

In the final part of the book I discuss how to create a system tailored for each of these
audiences. Before then, each time you see one or more of these heading boxes it indicates
that the material in that section of the book is aimed mainly at the relevant group and is
optional for others.

\subsection*{Asset allocating investor}

An asset allocating investor allocates funds amongst, and within, different asset classes.
Asset allocators can use systematic methods to avoid the short-term chasing of fads and
fashions that they know will reduce their returns. They might be lazy and wise amateur
investors, or managing institutional portfolios with long horizons such as pension funds.
Asset allocators are sceptical about those who claim to get extra returns from frequent
trading. For this reason the basic asset allocation example assumes you can't forecast how
asset prices will perform. However some investors might want to incorporate their views,
or the views of others. I show you how to achieve this without overtrading or ending up
with an extreme portfolio.

Unlike the other examples asset allocators usually don't use leverage. I illustrate the
investment process with the use of unleveraged passive exchange traded funds (ETFs).
But the methods I show apply equally well to investors in collective funds of both the
active and passive varieties, and to those investing in portfolios of individual securities.
My own portfolio includes a basket of ETFs which I manage using the principles of the
asset allocating investor.

\subsection*{Semi-automatic trader}

Semi-automatic traders live in a world of opportunistic bets\footnote{I am not using the
gambling term ``bet'' here in a pejorative sense. In my opinion the distinction some
people draw between financial gambling, trading and investing is completely meaningless:
they all involve taking financial risk on uncertain outcomes. Indeed professional gamblers
usually have a better understanding of risk management than many people working in the
investment industry.} taken on a fluid set of assets. Semi-automatic traders think they
are superior to simple rules when it comes to forecasting by how much prices will go up
or down; instead they make their own educated guesses. However they would like to
place those bets inside a systematic framework which will ensure their positions and risk
are properly managed. This frees them up to spend more time making the right call on the
market.

In my example the semi-automatic trader is comfortable with leverage and investing with
derivatives. They are both buyers and sellers, betting for or against asset prices. My
semi-automatic trader is active in equity index and commodity spread bet markets, but
the example is widely applicable elsewhere.

I trade a portfolio of UK equities using the framework I've outlined here for
semi-automatic traders.

\subsection*{Staunch systems trader}

The staunch systems trader is a true believer in the benefits of fully systematic trading.
Unlike the semi-automatic trader and the asset allocating investor, they embrace the use
of systematic trading rules to forecast price changes, but within the same common
framework for position risk management.

Many systems traders think they can find trading rules that give them extra profits, or
alpha. Others are unconvinced they have any special skill but believe there are additional
returns available which can't be captured just by ``buy and hold'' investing. They can use
very simple rules to capture these sources of alternative beta.

Systems traders may have access to back-testing software, either in off-the-shelf
packages, spreadsheets or bespoke software. It isn't absolutely necessary to have such
programs as I will be providing a flexible pre-configured trading system which doesn't
need back-testing. However if you want to develop your own new ideas I will show you
how to use these powerful software tools safely.

Like semi-automatic traders, staunch systems traders are comfortable with derivatives and
leverage. Although the examples I give are for futures trading, they are equally valid for
trading similar assets.

I trade over 40 futures contracts with my own money using a fully systematic set of eight
trading rules.

\section*{The technical stuff}
\phantomsection
\addcontentsline{toc}{section}{The technical stuff}

Inevitably this is a subject which requires some specialised terminology. Although I try
and keep jargon to a minimum it's usually easier to use a well-known shorthand term
rather than spelling everything out. Words and phrases highlighted in bold are briefly
defined in the glossary.

As well as standard finance vocabulary I use my own invented terms. So phrases like
instrument block are also in bold type, and appear in the glossary with a short
explanation.

Certain important concepts require deeper understanding and I will include more detail
when I first use them, as in the box below.

All detailed explanations in the text are signposted from the glossary, to help you refer to
them later in the book.

\conceptbox{%
\textbf{CONCEPT: EXAMPLE}

\medskip
These concept boxes give detailed explanations of key concepts in the book.
}

\section*{What is coming}
\phantomsection
\addcontentsline{toc}{section}{What is coming}

Running a systematic strategy is just like following any list of instructions, such as a
recipe. Many books on trading are like fast food outlets which give you something quick
and convenient to eat from a limited menu. However I am going to help you create your
own strategies; this is more than just a cookbook, it's a guide to writing recipes from
scratch. Inventing your own system requires extra work upfront, but it is more satisfying
and profitable in the long run.

To create your own recipes requires an understanding of the science of food chemistry and
of different kinds of dishes. Similarly, part one of this book --- ``Theory'' --- provides a
theoretical basis for why you should run systematic strategies and gives an overview of the
trading styles that are available.

Part two --- ``Toolbox'' --- provides you with two key methods used in the creation of
systematic strategies: back-testing and portfolio optimisation. Like sharp kitchen knives
these are powerful tools, but also potentially dangerous. When misused large trading
losses can be made despite apparently promising ideas. I will show you how to use them
properly --- and when you don't need them.

Part three --- ``Framework'' --- provides a complete and extendable framework for the
creation of systematic strategies. Finally in part four --- ``Practice'' --- I show three
different uses of the framework in action. I illustrate how it can be adapted for each of
the semi-automatic trader, asset allocating investor and staunch systems trader.

This book couldn't possibly be a comprehensive guide to the entire subject of trading.
Appendix A contains some books I would recommend for further reading and others I've
referred to in the text. There is also advice on where you can get data and access to
suitable brokers to begin investing systematically.

I've avoided putting mathematical formulas and detailed algorithms into the main text.
Instead appendix B contains detailed specifications for the trading rules, appendix C
covers portfolio optimisation, and appendix D includes further details on implementing
the framework. The website for this book --- \url{www.systematictrading.org} --- includes
additional material.

This is not a book about automating trading strategies. It's possible to trade
systematically using an entirely manual process with just a spreadsheet to speed up
calculations, so automation is not necessary. Nevertheless automation is desirable when
running fast and complex strategies. My website also includes some details on my own
automated system and guidance to help you develop your own.
